\section{Future Directions}
\label{sec:future}

Three directions follow directly from the map. \textbf{(1) Standardised evaluation of
self-improvement.} The systems on the top rung (\S3.1.5) each report improvement on their
own splits; the field lacks a shared protocol that measures held-out success as a function
of accumulated experience (an area-under-the-learning-curve metric), cross-embodiment
transfer, and library-reuse rate under a common budget. Building such a benchmark is a natural next step. \textbf{(2) Adaptation as a first-class marketplace
primitive.} The open problems of \S\ref{sec:open} suggest a research programme in which a
downloaded skill is not a frozen artefact but a starting point that the \S\ref{sec:code}
loop specialises to the local robot and environment; turning distribution (layer~2) and
adaptation (layer~3) into a single pipeline. \textbf{(3) Bridging the RL and code flavours
of skills.} \S\ref{sec:skill} shows that skills are learned either as latent RL policies or
as retrievable code; a unifying interface, code that wraps, calls, and refines learned
policies, and learned policies distilled back into inspectable code, would let the
self-improvement machinery of \S\ref{sec:code} operate over both. We see the combination of
a standard self-improvement benchmark with an adaptation-centric marketplace as the most
likely path from today's static skill stores to genuinely capable, compounding robots.
